\chapter[Setting: Towers, Kronecker Map, Folding]{Setting: Polynomial Towers, Kronecker Embedding and Symbolic Folding}
\label{ch:setting}

This chapter fixes all objects used in the rest of the text: the fields, the evaluation
domains and their fibres, the Reed--Solomon codes, the tower Kronecker embedding, the
fold operator, and the randomized symbolic FRI protocol. It ends with the honest
invariant, which identifies the terminal polynomial $g(T)$.

The domains are \emph{polynomial towers}: a finite set $D_0$ and polynomial maps
$\sigma_1,\dots,\sigma_k$ whose restrictions to successive images have fibres of constant
size. Multiplicative cosets with power maps (smooth FRI domains) and affine subspaces
of binary fields with subspace polynomials (additive FRI domains) are the two main
instances (\cref{ex:multiplicative,ex:additive}).

%=====================================================================
\section{Fields and polynomial towers}
%=====================================================================

\begin{definition}[Fields]\label{def:fields}
Let $q$ be a prime power, $\F:=\F_q$ the \emph{base field}, $e\ge 1$ an integer, and
$\K:=\F_{q^e}$ the \emph{challenge field}. All verifier challenges, all coefficients in the
formal variable $T$, and all codes of \cref{def:rs} are over $\K$; the initial committed
word is over $\F$.
\end{definition}

\begin{definition}[Polynomial tower]\label{def:domain}
Fix integers $k\ge1$ and $a_1,\dots,a_k\ge2$ (the \emph{arities}), and set
\[
A_0:=1,\qquad A_i:=a_1a_2\cdots a_i\quad(1\le i\le k),\qquad L:=A_k .
\]
Fix an integer $B\ge2$ (the \emph{blowup}), $N:=BL$ and $\rho:=1/B$ (the \emph{rate}).
A \emph{polynomial tower} is a set $D_0\subseteq\F$ with $|D_0|=N$ together with
polynomials $\sigma_1,\dots,\sigma_k\in\F[X]$, $\deg\sigma_i=a_i$, such that, with
\[
D_i:=\sigma_i(D_{i-1})\quad(1\le i\le k),
\]
every $y\in D_{i+1}$ has exactly $a_{i+1}$ preimages under $\sigma_{i+1}$ in $D_i$
($0\le i<k$). Write $n_i:=N/A_i$, $L_i:=L/A_i$, and for $0\le i\le l\le k$
\[
\Sigma_{i,l}:=\sigma_l\circ\sigma_{l-1}\circ\cdots\circ\sigma_{i+1}\in\F[X],\qquad
\Sigma_{i,i}:=X,\qquad \Sigma_l:=\Sigma_{0,l}.
\]
\end{definition}

\begin{lemma}[Fibres]\label{lem:fibres}
For $0\le i<k$ and $y\in D_{i+1}$ put $\fib(y):=\{x\in D_i:\sigma_{i+1}(x)=y\}$. Then
$|\fib(y)|=a_{i+1}$, $D_i=\bigsqcup_{y\in D_{i+1}}\fib(y)$, $|D_i|=n_i$, $n_i=a_{i+1}n_{i+1}$,
$n_k=B$, $L_i/n_i=\rho$, and $\deg\Sigma_{i,l}=A_l/A_i$.
\end{lemma}

\begin{proof}
The fibres partition $D_i$ and each has $a_{i+1}$ elements by assumption, so
$|D_i|=a_{i+1}|D_{i+1}|$; induction from $|D_0|=N$ gives $|D_i|=N/A_i=n_i$. The remaining
identities follow from $L=A_k$, $N=BL$ and multiplicativity of degrees under composition.
\end{proof}

\begin{example}[Multiplicative towers]\label{ex:multiplicative}
Assume $N\mid q-1$, let $H\le\F^\times$ be the subgroup of order $N$, $x_0\in\F^\times$,
$D_0:=x_0H$ and $\sigma_i(X):=X^{a_i}$. Then $D_i=x_0^{A_i}H^{A_i}$, and for $y\in D_{i+1}$
and any $x\in\fib(y)$, $\fib(y)=x\,\mu_{a_{i+1}}$, where $\mu_\ell$ is the group of
$\ell$-th roots of unity: $H^{A_i}$ is cyclic of order $n_i$ and $a_{i+1}\mid n_i$, so the
kernel of $z\mapsto z^{a_{i+1}}$ on $H^{A_i}$ has order $a_{i+1}$. Here
$\Sigma_{i,l}(X)=X^{A_l/A_i}$. This is the smooth FRI domain; the binary case is $a_i=2$.
\end{example}

\begin{example}[Additive towers]\label{ex:additive}
Let $q=2^t$, let $U_0\subseteq\F$ be an $\F_2$-subspace of dimension $k+\log_2B$ (with
$B$ a power of $2$), $D_0:=x_0+U_0$, and let $\sigma_{i}$ be $\F_2$-linearized polynomials
of degree $2$ with $\ker\sigma_{i}\subseteq\Sigma_{i-1}(U_0)$ of order $2$ (for instance
the normalized subspace polynomials of a flag of $U_0$, as in~\cite{LCH14,DP24}). Then
each $\sigma_i$ is $2$-to-$1$ on the affine subspace $D_{i-1}$, with fibres the cosets of
$\ker\sigma_i$, and \cref{def:domain} holds with $a_i=2$. Higher arities arise from
kernels of order $2^b$.
\end{example}

%=====================================================================
\section{Reed--Solomon codes and agreement}
%=====================================================================

\begin{definition}[Reed--Solomon codes]\label{def:rs}
For $0\le i\le k$,
\[
\RS_i:=\bigl\{\,G|_{D_i}\;:\;G\in\K[X],\ \deg G<L_i\,\bigr\}\subseteq \K^{D_i}.
\]
Its rate is $L_i/n_i=\rho$ and its minimum distance is $n_i-L_i+1$. Since $L_i\le n_i$,
the map $G\mapsto G|_{D_i}$ is injective on $\K[X]_{<L_i}$, and we identify a codeword with
its polynomial.
\end{definition}

\begin{definition}[Pointwise and block agreement]\label{def:agreement}
For words $w,v:D_i\to\K$ (or $\K[T]$), the \emph{relative agreement} is
$\mathrm{agr}(w,v):=\frac1{n_i}\#\{x\in D_i: w(x)=v(x)\}$ and $\Delta:=1-\mathrm{agr}$.
For $i<k$, the \emph{block agreement} is
\[
\agr(w,v):=\frac1{n_{i+1}}\#\bigl\{y\in D_{i+1}:\ w|_{\fib(y)}=v|_{\fib(y)}\bigr\},
\qquad \Db:=1-\agr .
\]
For a set $\mathcal C$ of words, $\agr(w,\mathcal C):=\max_{v\in\mathcal C}\agr(w,v)$.
\end{definition}

\begin{lemma}[Fibre agreement of distinct codewords]\label{lem:codeword-fibres}
If $G\neq G'$ with $\deg G,\deg G'<L_i$ and $i<k$, then
\[
\#\{y\in D_{i+1}: G|_{\fib(y)}=G'|_{\fib(y)}\}\le \frac{L_i-1}{a_{i+1}}<L_{i+1}=\rho\,n_{i+1}.
\]
\end{lemma}

\begin{proof}
Each fibre on which $G$ and $G'$ agree contributes $a_{i+1}$ distinct roots of the nonzero
polynomial $G-G'$, which has at most $L_i-1$ roots; and $L_i=a_{i+1}L_{i+1}$.
\end{proof}

%=====================================================================
\section{Tower expansions and the Kronecker embedding}
%=====================================================================

\begin{lemma}[$\sigma$-adic expansion]\label{lem:sigma-adic}
Let $R$ be a field or $R=\K[T]$, let $\sigma\in\F[X]$ have degree $a\ge1$, and let $M\ge1$.
Every $P\in R[X]$ with $\deg_XP<aM$ has a unique expansion
\[
P(X)=\sum_{j<a}X^j\,P_j\bigl(\sigma(X)\bigr),\qquad P_j\in R[X],\ \deg_XP_j<M .
\]
\end{lemma}

\begin{proof}
It suffices to treat $R=\K$; for $R=\K[T]$ apply the statement coefficientwise in $T$.
The map $(P_j)_{j<a}\mapsto\sum_jX^jP_j(\sigma)$ is $\K$-linear between spaces of dimension
$aM$. If $P_j\ne0$ then $\deg\bigl(X^jP_j(\sigma)\bigr)=j+a\deg P_j$, and these degrees are
pairwise distinct for distinct $j$ (they differ modulo $a$). Hence a nontrivial sum has
degree equal to the largest of them and is nonzero: the map is injective, hence bijective.
\end{proof}

\begin{definition}[Degree-bounded polynomials]\label{def:Pa}
For a field $E\in\{\F,\K\}$,
\[
\Pa(E):=\bigl\{f\in E[Y_1,\dots,Y_k]:\ \deg_{Y_i}f<a_i\ \text{for }1\le i\le k\bigr\}.
\]
For $\mathbf a=(2,\dots,2)$ this is the space of multilinear polynomials.
\end{definition}

\begin{definition}[Tower Kronecker map]\label{def:kronecker}
The map $\kappa:\Pa(E)\to E[X]_{<L}$ is
\[
\kappa(f)(X):=f\bigl(\Sigma_0(X),\Sigma_1(X),\dots,\Sigma_{k-1}(X)\bigr).
\]
For $G\in E[X]_{<L}$ write $f_G:=\kappa^{-1}(G)$ (well defined by \cref{lem:kronecker}).
\end{definition}

\begin{lemma}\label{lem:kronecker}
$\kappa$ is an $E$-linear bijection. For multiplicative towers
(\cref{ex:multiplicative}) the monomial $Y_1^{e_1}\cdots Y_k^{e_k}$ is sent to
$X^{\sum_ie_iA_{i-1}}$ (mixed-radix Kronecker substitution); for additive towers with
normalized subspace polynomials (\cref{ex:additive}) the monomials $\prod_{i\in S}Y_i$ are
sent to the novel polynomial basis $\prod_{i\in S}\Sigma_{i-1}(X)$ of~\cite{LCH14}.
\end{lemma}

\begin{proof}
A monomial $Y^{\mathbf e}$ with $e_i<a_i$ is sent to $\prod_i\Sigma_{i-1}^{e_i}$, of degree
$\sum_ie_iA_{i-1}\le\sum_i(a_i-1)A_{i-1}=L-1$ (telescoping), so $\kappa$ lands in
$E[X]_{<L}$; both spaces have dimension $L$. Injectivity: by induction on $k$. For $k=1$,
$\kappa(f)=f(X)$. For $k\ge2$, write $f=\sum_{j<a_1}Y_1^jf_j(Y_2,\dots,Y_k)$; then
\[
\kappa(f)(X)=\sum_{j<a_1}X^j\,\kappa'(f_j)\bigl(\sigma_1(X)\bigr),
\]
where $\kappa'$ is the map of the tower $(\sigma_2,\dots,\sigma_k)$ (because
$\Sigma_{l}=\Sigma_{1,l}\circ\sigma_1$), and $\deg\kappa'(f_j)<L/a_1$. If $\kappa(f)=0$, then
$\kappa'(f_j)=0$ for all $j$ by uniqueness in \cref{lem:sigma-adic}, hence $f_j=0$ by
induction. The two explicit descriptions are the cases $\Sigma_{i-1}=X^{A_{i-1}}$ and
$\Sigma_{i-1}=$ the normalized subspace polynomial of dimension $i-1$.
\end{proof}

%=====================================================================
\section{Components and the fold operator}
%=====================================================================

Throughout this section fix $0\le i<k$ and write $a:=a_{i+1}$, $\sigma:=\sigma_{i+1}$. Let
$R$ denote either $\K$ or $\K[T]$; words take values in $R$.

\begin{definition}[Components]\label{def:components}
For $y\in D_{i+1}$ let $M_y:=\bigl(x^j\bigr)_{x\in\fib(y),\,0\le j<a}$, an $a\times a$
matrix (rows indexed by $\fib(y)$ in a fixed order). For $w:D_i\to R$ define
$u[w](y)=(u_j[w](y))_{j<a}\in R^a$ by
\[
u[w](y):=M_y^{-1}\,\bigl(w(x)\bigr)_{x\in\fib(y)} .
\]
\end{definition}

\begin{lemma}[Components]\label{lem:components}
\begin{enumerate}[label=\textup{(\roman*)}]
\item $M_y$ is invertible, and $u[w](y)$ does not depend on the order of $\fib(y)$.
\item For every $y\in D_{i+1}$ the map
$R^{\fib(y)}\to R^a$, $w|_{\fib(y)}\mapsto u[w](y)$, is an $R$-linear bijection, with
inverse $w(x)=\sum_{j<a}x^ju_j[w](\sigma(x))$.
\item If $w=G|_{D_i}$ with $G=\sum_{j<a}X^jG_j(\sigma(X))$, $G_j\in R[X]$, then
$u_j[w]=G_j|_{D_{i+1}}$. In particular, if $\deg_XG<L_i$ then $u_j[w]\in\RS_{i+1}$
\textup{(}coefficientwise in $T$ when $R=\K[T]$\textup{)}.
\end{enumerate}
\end{lemma}

\begin{proof}
(i) $M_y$ is a Vandermonde matrix on the $a$ distinct points of $\fib(y)$; reordering the
rows reorders both sides of $M_yu=(w(x))_x$ consistently.
(ii) By definition $(w(x))_x=M_y\,u[w](y)$, which is the inversion formula; $M_y$ is
invertible.
(iii) For $x\in\fib(y)$, $G(x)=\sum_jx^jG_j(y)$, i.e.\ $(G(x))_x=M_y(G_j(y))_j$, so
$u_j[w](y)=G_j(y)$ by (ii). If $\deg_XG<L_i=aL_{i+1}$ then $\deg G_j<L_{i+1}$ by
\cref{lem:sigma-adic}.
\end{proof}

\begin{remark}[Multiplicative components]\label{rem:dft}
For $\sigma=X^a$ and $\fib(y)=x\mu_a$, orthogonality of characters gives the explicit
inverse $u_j[w](y)=\frac1a\sum_{\omega\in\mu_a}(\omega x)^{-j}w(\omega x)$; for $a=2$,
$u_0=\frac{w(x)+w(-x)}2$ and $u_1=\frac{w(x)-w(-x)}{2x}$.
\end{remark}

\begin{definition}[Fold]\label{def:fold}
For $w:D_i\to R$ and $\beta\in R$,
\[
\Fold_\beta[w]:D_{i+1}\to R,\qquad \Fold_\beta[w]:=\sum_{j<a}\beta^j\,u_j[w].
\]
\end{definition}

\begin{lemma}[Fold of a polynomial]\label{lem:fold-poly}
If $G=\sum_{j<a}X^jG_j(\sigma(X))\in R[X]$ with $\deg_XG<L_i$, then
$\Fold_\beta[G|_{D_i}]=\bigl(\sum_{j<a}\beta^jG_j\bigr)\big|_{D_{i+1}}$, a polynomial of
$X$-degree $<L_{i+1}$. In particular $\Fold_\beta$ maps $\RS_i$ into $\RS_{i+1}$ for
$\beta\in\K$.
\end{lemma}

\begin{proof}
\Cref{lem:components}(iii) and linearity.
\end{proof}

\begin{definition}[Full fold]\label{def:full-fold}
For $G\in\K[X]_{<L}$ and $\beta=(\beta_1,\dots,\beta_k)\in\K^k$ put
$G^{(0)}:=G$ and $G^{(i+1)}:=\Fold_{\beta_{i+1}}[G^{(i)}]$ (as polynomials, via
\cref{lem:fold-poly}). Then $G^{(k)}$ has degree $<L_k=1$, i.e.\ is a constant, denoted
$\widehat G(\beta)$.
\end{definition}

\begin{lemma}[Full fold evaluates $f_G$]\label{lem:full-fold}
For all $G\in\K[X]_{<L}$ and $\beta\in\K^k$: $\widehat G(\beta)=f_G(\beta_1,\dots,\beta_k)$.
\end{lemma}

\begin{proof}
Apply \cref{prop:honest} below with $f=f_G$ and $r=\beta$: by part (ii) and
\cref{lem:specialization} with $c=1$, the scalar fold chain of $G$ with challenges $\beta$
is $F_i(\,\cdot\,,1)$, and by part (iii) its last term is $g(1)=f_G(\beta)$.
\end{proof}

%=====================================================================
\section{The randomized symbolic FRI protocol}
%=====================================================================

\begin{definition}[Degree caps]\label{def:degree-caps}
$d_0:=0$, $d_i:=\sum_{l=1}^{i}(a_l-1)$ for $1\le i\le k$, and $d:=d_k$.
For $\mathbf a=(2,\dots,2)$, $d_i=i$ and $d=k$.
\end{definition}

\begin{protocol}[Randomized symbolic FRI]\label{prot:sfri}
Public parameters: $\F,\K$, a polynomial tower $(D_0,\sigma_1,\dots,\sigma_k)$, and a number
$s\ge1$ of query paths. Words at layer $i\ge1$ take values in $\K[T]_{\le d_i}$, encoded as
fixed-length coefficient vectors; leaves are grouped by fibre.

\smallskip\noindent\textbf{Commit phase.}
\begin{enumerate}[label=\textup{C\arabic*.}]
\item $\mathsf P$ sends an oracle $W_0:D_0\to\F$.
\item For $i=0,\dots,k-1$: $\mathsf V$ samples $r_{i+1}\leftarrow\K$ uniformly and sends it;
$\mathsf P$ sends an oracle $W_{i+1}:D_{i+1}\to\K[T]_{\le d_{i+1}}$.
\item $\mathsf P$ sends $\gt\in\K[T]$; $\mathsf V$ rejects unless $\deg\gt\le d$.
\end{enumerate}

\smallskip\noindent\textbf{Query phase.} Repeat $s$ times independently:
\begin{enumerate}[label=\textup{Q\arabic*.}]
\item $\mathsf V$ samples $x_0\leftarrow D_0$ uniformly and sets
$x_{i+1}:=\sigma_{i+1}(x_i)$ for $0\le i<k$.
\item For $0\le i<k$, $\mathsf V$ reads $W_i$ on $\fib(x_{i+1})$ and $W_{i+1}(x_{i+1})$ and
checks the identity in $\K[T]$
\[
W_{i+1}(x_{i+1})=\Fold_{r_{i+1}T}[W_i](x_{i+1}).
\]
\item $\mathsf V$ checks $W_k(x_k)=\gt$ in $\K[T]$.
\end{enumerate}
$\mathsf V$ accepts iff all checks of all $s$ repetitions pass. The event of acceptance is
denoted $\Acc$.
\end{protocol}

\begin{remark}[Variants covered by the analysis]\label{rem:variants}
\begin{enumerate}[label=\textup{(\roman*)}]
\item The oracle $W_k$ may be omitted, in which case Q3 is merged into Q2 for $i=k-1$ as
$\Fold_{r_kT}[W_{k-1}](x_k)=\gt$. Alternatively $\mathsf V$ may read all $B$ values of $W_k$.
All results of this text hold for each variant.
\item The values of $W_0$ lie in $\F\subset\K[T]_{\le 0}$ by the encoding; a leaf that is not
a constant is rejected. This is the only degree condition used by the soundness analysis
besides $\deg\gt\le d$; the caps $d_i$ for $1\le i<k$ only fix the encoding length.
\item Since all fibres have the same size, each $x_i$ in Q1 is uniform on $D_i$.
\item The challenges $r_{i+1}T$ may be replaced by affine or polynomial curves in $T$;
see \cref{sec:curves}.
\end{enumerate}
\end{remark}

%=====================================================================
\section{Honest invariant and completeness}
%=====================================================================

\begin{proposition}[Honest invariant]\label{prop:honest}
Let $f\in\Pa(\K)$ and $r=(r_1,\dots,r_k)\in\K^k$. For $0\le i\le k$ define
\[
F_i(X,T):=f\Bigl(r_1T,\dots,r_iT,\;\Sigma_{i,i}(X),\;\Sigma_{i,i+1}(X),\;\dots,\;\Sigma_{i,k-1}(X)\Bigr)
\in\K[X,T].
\]
Then
\begin{enumerate}[label=\textup{(\roman*)}]
\item $F_0=\kappa(f)$, $\deg_XF_i<L_i$ and $\deg_TF_i\le d_i$;
\item $F_{i+1}|_{D_{i+1}}=\Fold_{r_{i+1}T}\bigl[F_i|_{D_i}\bigr]$ for $0\le i<k$;
\item $F_k=g(T):=f(r_1T,\dots,r_kT)$, and $\deg g\le d$.
\end{enumerate}
\end{proposition}

\begin{proof}
(i) $F_0=\kappa(f)$ by definition. In $F_i$ the variable $Y_l$ ($l>i$) is replaced by
$\Sigma_{i,l-1}(X)$, of degree $A_{l-1}/A_i$, so a monomial $Y^{\mathbf e}$ contributes
$X$-degree $\sum_{l>i}e_lA_{l-1}/A_i\le\sum_{l>i}(a_l-1)A_{l-1}/A_i=L_i-1$, and $T$-degree
$\sum_{l\le i}e_l\le d_i$.

(ii) Write $f=\sum_{j<a_{i+1}}Y_{i+1}^j f_j$ with $f_j$ free of $Y_{i+1}$. Since
$\Sigma_{i,i}(X)=X$ and $\Sigma_{i,l}=\Sigma_{i+1,l}\circ\sigma_{i+1}$ for $l\ge i+1$,
\begin{align*}
F_i(X,T)&=\sum_{j<a_{i+1}}X^j\,\phi_j\bigl(\sigma_{i+1}(X),T\bigr),\\
\phi_j(Z,T)&:=f_j\bigl(r_1T,\dots,r_iT,\,\cdot\,,\Sigma_{i+1,i+1}(Z),\dots,\Sigma_{i+1,k-1}(Z)\bigr),
\end{align*}
where the slot of $Y_{i+1}$ is empty since $f_j$ does not depend on it. By (i) applied to
$f_j$, $\deg_Z\phi_j<L_{i+1}$. By \cref{lem:fold-poly} with $R=\K[T]$ and
$\beta=r_{i+1}T$,
\[
\Fold_{r_{i+1}T}[F_i|_{D_i}]=\Bigl(\sum_{j}(r_{i+1}T)^j\phi_j(Y,T)\Bigr)\Big|_{Y\in D_{i+1}} .
\]
The polynomial $\sum_j(r_{i+1}T)^j\phi_j(Y,T)$ is $f$ with $Y_l=r_lT$ for $l\le i+1$ and
$Y_l=\Sigma_{i+1,l-1}(Y)$ for $l\ge i+2$, i.e.\ it is $F_{i+1}(Y,T)$.
(iii) For $i=k$ no variable remains; $\deg g\le d_k=d$ by (i).
\end{proof}

\begin{corollary}[Completeness]\label{cor:completeness}
If $W_0=\kappa(f)|_{D_0}$ for some $f\in\Pa(\F)$, $W_i=F_i|_{D_i}$ for $1\le i\le k$ and
$\gt=g$, then $\mathsf V$ accepts with probability $1$. Moreover $\deg_TW_i\le d_i$, so the
encoding of \cref{prot:sfri} suffices.
\end{corollary}

\begin{proof}
\Cref{prop:honest}(ii),(iii); $F_0$ has coefficients in $\F$ because $f$ and the $\sigma_i$ do.
\end{proof}

\begin{lemma}[Scalar specialization]\label{lem:specialization}
Let $w:D_i\to\K[T]$, $\varphi\in\K[T]$ and $c\in\K$. Writing $w^{(c)}(x):=w(x)(c)$,
\[
\Fold_{\varphi}[w](y)\big|_{T=c}=\Fold_{\varphi(c)}\bigl[w^{(c)}\bigr](y)\qquad(y\in D_{i+1}).
\]
In particular $\Fold_{rT}[w](y)|_{T=c}=\Fold_{cr}[w^{(c)}](y)$, and every check Q2/Q3 that
holds in $\K[T]$ holds after evaluation at $T=c$ for the scalar challenge $cr_{i+1}$.
\end{lemma}

\begin{proof}
$u_j[w](y)$ is a $\K$-linear combination of the values $w(x)$, $x\in\fib(y)$, with
coefficients the entries of $M_y^{-1}$; so evaluation at $T=c$ commutes with $u_j$, and
$\varphi^j|_{T=c}=\varphi(c)^j$.
\end{proof}

\begin{remark}[Binary multiplicative case]\label{rem:binary}
For the multiplicative tower with $\mathbf a=(2,\dots,2)$, $k=m$: $\Pa$ is the
multilinear space, $\kappa$ is the Kronecker substitution $X_i\mapsto X^{2^{i-1}}$, $d=m$,
and \cref{prot:sfri} is the protocol of the original handoff, with
$g(T)=f(r_1T,\dots,r_mT)$.
\end{remark}
